\subsection{复合算子的指示函数}

在与 \(n^{\circ}\) 5 相同的假设下，把（VI, p.~277）的公式 (26) 应用于函数 \(e^{\lambda x}\) ，就得到

\[
U _ {x} ^ {\xi} (e ^ {\lambda \xi}) = U _ {0} ^ {\xi} (e ^ {\lambda x} e ^ {\lambda \xi}) = e ^ {\lambda x} U _ {0} ^ {\xi} (e ^ {\lambda \xi}) = u (\lambda) e ^ {\lambda x}\tag{30}
\]

其中令 \(u(\lambda) = U_{0}^{\xi}(e^{\lambda\xi})\) 。定义在 C 上、取复数值的函数 \(u(\lambda)\) 称为复合算子 U 的指示函数。注意，若 U 在多项式环 C{[}X{]} 上的限制等于级数

\[
\mathbf {D} ^ {p} \sum_ {n = 0} ^ {\infty} \alpha_ {n} \mathbf {D} ^ {n}
\]

\pandocbounded{\includegraphics[keepaspectratio]{images/part002_816cd56c.jpg}}

（VI, p.~271, th. 1）（我们在 VI, p.~271 已把它记作 \(u(\mathrm{D})\) ），则以 \(\alpha_{n}\lambda^{n+p}\) 为通项的复数项级数对 \(\lambda \neq 0\) 不一定收敛，而且即使它对 \(\lambda\) 的某些值收敛，其和也不一定等于 U 的指示函数 \(u(\lambda)\) （VI, p.~291, exerc. 2）。若存在 C 中 0 的一个邻域，使得以 \(\alpha_{n}\lambda^{n+p}\) 为通项的级数在其上绝对收敛且其和等于指示函数 \(u(\lambda)\) ，则我们称复合算子 U 为正则的 \(^{1}\) 。把 VI, p.~278 的公式 (27) 应用于函数 \(e^{\lambda x}\) ，取 h = 0；由于 \(\mathrm{D}^{m}(e^{\lambda x}) = \lambda^{m} e^{\lambda x}\) ，有 \(U_{0}^{\xi}\left(\mathrm{D}^{m}(e^{\lambda x})\right) = \lambda^{m} u(\lambda)\) ；由此得出，对每个复数 \(\lambda\) （满足 \(u(\lambda) \neq 0\) 的）

\[
\frac {\lambda^ {p} e ^ {\lambda x}}{u (\lambda)} = \sum_ {m = 0} ^ {n} u _ {m} (x) \frac {\lambda^ {m}}{m !} - \frac {\lambda^ {n + 1}}{u (\lambda)} U _ {0} ^ {\xi} \left(\int_ {0} ^ {\xi - x} \frac {1}{n !} u _ {n} (x + \eta) e ^ {\lambda (\xi - \eta)} d \eta\right)\tag{31}
\]

特别地，对 \(x = 0\) ，有

\[
\frac {\lambda^ {p}}{u (\lambda)} = \sum_ {m = 0} ^ {n} \beta_ {m} \frac {\lambda^ {m}}{m !} - \frac {\lambda^ {n + 1}}{u (\lambda)} U _ {0} ^ {\xi} \left(\int_ {0} ^ {\xi} \frac {1}{n !} u _ {n} (\eta) e ^ {\lambda (\xi - \eta)} d \eta\right)\tag{32}
\]

其中 \(\beta_{m} = u_{m}(0)\) 。

若 \(U\) 是正则算子，则对一切 \(\lambda \in \mathbf{C}\) （使幂级数 \(u(\lambda) = \sum_{n=0}^{\infty} \alpha_n \lambda^{n+p}\) 与 \(\sum_{n=0}^{\infty} \beta_n \frac{\lambda^n}{n!}\) 绝对收敛者）\(^2\) ，由公式 (16) 及两个绝对收敛级数的乘积公式（Gen.~Top., VIII, p.~115, prop. 1），有

\[
{\frac {\lambda^ {p}}{u (\lambda)}} = \sum_ {n = 0} ^ {\infty} \beta_ {n} {\frac {\lambda^ {n}}{n !}}.\tag{33}
\]

同样，由于 \(e^{\lambda x}\) 的泰勒展开对每个 \(\lambda \in \mathbf{C}\) 和每个 \(x \in \mathbf{C}\) 都绝对收敛（III, p.~106），（由公式 (6) （VI, p.~273）与 (16) （VI, p.~274））对所考虑的一切值和一切 \(x \in \mathbf{C}\) 也有

\[
\frac {\lambda^ {p} e ^ {\lambda x}}{u (\lambda)} = \sum_ {n = 0} ^ {\infty} u _ {n} (x) \frac {\lambda^ {n}}{n !}.\tag{34}
\]

注. 可以利用公式 (33) （相应地 (34)），借助下面关于幂级数的引理来计算诸 \(\beta_{n}\) （相应地诸 \(u_{n}(x)\) ）：

引理. 若两个幂级数 \(\sum_{n=0}^{\infty} c_n \lambda^n\) 、\(\sum_{n=0}^{\infty} d_n \lambda^n\) 在 0 的某个邻域内的一切 \(\lambda\) 处绝对收敛，且若 \(\sum_{n=0}^{\infty} c_n \lambda^n = \sum_{n=0}^{\infty} d_n \lambda^n\) 对这些 \(\lambda\) 的值成立，则 \(c_n = d_n\) （对一切整数 \(n \geqslant 0^3\) ）。

若用任何方法能找到一个在 0 的某个邻域上收敛于 \(\lambda^p / u(\lambda)\) 的幂级数，则该级数的系数必然等于诸 \(\beta_n / n!\) 。在下面各例中我们就要应用这个方法。

例. 1) 若 \(U\) 是恒等映射，则有 \(u(\lambda) = 1\) ，且显然算子 \(U\) 是正则的；由于 \(u_{n}(x) = x^{n}\) ，令 \(t = \xi - \eta\) ，VI, p.~278 的公式 (27) 可以写成

\[
f (x + h) = \sum_ {n = 0} ^ {\infty} \frac {1}{m !} f ^ {(m)} (h) x ^ {m} + \int_ {h} ^ {x + h} f ^ {(n + 1)} (t) \frac {(x + h - t) ^ {n}}{n !} d t
\]

也就是说，它归结为泰勒公式（II, p.~62）。

\begin{enumerate}
\def\labelenumi{\arabic{enumi})}
\setcounter{enumi}{1}
\tightlist
\item
  我们取 U 为这样的复合算子：它把定义在 \(R_{+}\) 上的每个函数 f 对应到函数 \(x \mapsto f(x + 1) - f(x)\) ；于是
\end{enumerate}

\[
U _ {x} ^ {\xi} (f (\xi)) = f (x + 1) - f (x);
\]

我们已看到（VI, p.~275），U 在 C{[}X{]} 上的限制等于 \(e^{D}-1\) 。另一方面，由于 \(u(\lambda)=e^{\lambda}-1\) ，算子 U 是正则的；我们将看到（VI, p.~288）如何确定伯努利数 \(b_{n}\) ，办法是把 \(\frac{\lambda}{e^{\lambda}-1}\) 展开为收敛的幂级数。把 VI, p.~278 的公式 (27) 应用于一个函数 f 的原函数（该函数有连续的 \(n^{th}\) 阶导数，在 \(R_{+}\) 上），就得到

\[
\begin{array}{l} f (x + h) = \int_ {h} ^ {h + 1} f (t) d t \\ \qquad + \sum_ {m = 1} ^ {n} \frac {1}{m !} \mathrm{B} _ {m} (x) \left(f ^ {(m - 1)} (h + 1) - f ^ {(m - 1)} (h)\right) + \mathrm{R} _ {n} (x, h) \end{array}\tag{35}
\]

\[
\begin{array}{l} \mathrm{R} _ {n} (x, h) = - \int_ {0} ^ {1 - x} \frac {\mathrm{B} _ {n} (x + \eta)}{n !} f ^ {(n)} (h + 1 - \eta) d \eta \\ \qquad \qquad + \int_ {0} ^ {- x} \frac {\mathrm{B} _ {n} (x + \eta)}{n !} f ^ {(n)} (h - \eta) d \eta . \end{array}\tag{36}
\]

*3) 设 \(\mathbf{E}\) 是由函数 \(f\) （在 \(\mathbf{R}\) 上定义且连续、并使积分 \(\int_{-\infty}^{+\infty} f(x + \xi)e^{-\xi^2 / 2}d\xi\) 对一切 \(x \geqslant 0\) 收敛的）所组成的向量空间。算子 \(\mathrm{U}\) 由

\[
U _ {x} ^ {\xi} (f (\xi)) = \frac {1}{\sqrt {2 \pi}} \int_ {- \infty} ^ {+ \infty} e ^ {- \xi^ {2} / 2} f (x + \xi) d \xi
\]

定义，它这时定义在 E 上且显然是一个复合算子。空间 E 包含所有指数函数 \(e^{\lambda x}\) （ \(\lambda\) 为任意复数），并且有

\[
\begin{array}{l} u (\lambda) = \frac {1}{\sqrt {2 \pi}} \int_ {- \infty} ^ {+ \infty} e ^ {- (\xi^ {2} / 2) + \lambda \xi} d \xi \\ = \frac {1}{\sqrt {2 \pi}} e ^ {\lambda^ {2} / 2} \int_ {- \infty} ^ {+ \infty} e ^ {- (\xi - \lambda) ^ {2} / 2} d \xi = e ^ {\lambda^ {2} / 2} \end{array}
\]

（参见 III, p.~120, exerc. 24 及 VII, p.~313, 公式 (22)）。有 \(n! \alpha_n = U_0^\xi (\xi^n) = \frac{1}{\sqrt{2\pi}} \int_{-\infty}^{+\infty} e^{-\xi^2 / 2} \xi^n d\xi\) 。对每个整数 \(n\) 可以写

\[
\sum_ {k = 0} ^ {n} \int_ {- \infty} ^ {+ \infty} \frac {| \lambda \xi | ^ {k}}{k !} e ^ {- \xi^ {2} / 2} d \xi \leqslant 2 \int_ {0} ^ {+ \infty} e ^ {- (\xi^ {2} / 2) + | \lambda | \xi} d \xi .
\]

于是级数 \(\sum_{n=0}^{\infty}e^{-\xi^{2}/2}\frac{(\lambda\xi)^{n}}{n!}\) 可以在 R 上逐项积分（II, p.~72, cor. 1），这就证明级数 \(\sum_{n=0}^{\infty}\alpha_{n}\lambda^{n}\) 对每个 \(\lambda\in C\) 绝对收敛，且其和等于 \(u(\lambda)=e^{\lambda^{2}/2}=\sum_{n=0}^{\infty}\frac{\lambda^{2n}}{2^{n}n!}\) ；因而算子 U 是正则的。应用上述引理可知，有 \(\alpha_{2n}=1/2^{n}n!\) ，\(\alpha_{2n+1}=0\) （对每个 \(n\geqslant0\) ）；于是算子 U 是 0 阶的。有

\[
\frac {1}{u (\lambda)} = e ^ {- \lambda^ {2} / 2} = \sum_ {n = 0} ^ {\infty} \frac {(- 1) ^ {n} \lambda^ {2 n}}{2 ^ {n} n !},
\]

该级数对每个 \(\lambda \in \mathbf{C}\) 绝对收敛；再次应用该引理可知 \(\beta_{2n} = \frac{(-1)^n}{2^n}\frac{(2n)!}{n!}\) ，\(\beta_{2n+1} = 0\) ；此外，级数 \(\sum_{n=0}^{\infty}\frac{\lambda^n}{n!} u_n(x)\) 对每个 \(\lambda \in \mathbf{C}\) 和每个 \(x \in \mathbf{R}\) 绝对收敛，并且有

\[
\sum_ {n = 0} ^ {\infty} \frac {\lambda^ {n}}{n !} u _ {n} (x) = \exp \left(- \frac {\lambda^ {2}}{2} + \lambda x\right) = \exp \left(\frac {x ^ {2}}{2}\right) \exp \left(- \frac {1}{2} (\lambda - x) ^ {2}\right).
\]

把泰勒公式应用于函数 \(\exp(-x^2/2)\) ，就得到多项式 \(u_n(x)\) 的下列表达式：

\[
u _ {n} (x) = (- 1) ^ {n} e ^ {x ^ {2} / 2} \frac {d ^ {n}}{d x ^ {n}} (e ^ {- x ^ {2} / 2}).
\]

这个多项式称为 n 次埃尔米特多项式，常记作 \(\mathrm{H}_{n}(x)\) 。VI, p.~273 的公式 (7)、(8) 和 (9) 在这里给出

\[
\frac {d \mathrm{H} _ {n}}{d x} = n \mathrm{H} _ {n - 1} (x)
\]

\[
\mathrm{H} _ {n} (x + y) = \sum_ {k = 0} ^ {n} {\binom {n} {k}} \mathrm{H} _ {n - k} (x)   y ^ {k}
\]

\[
\frac {1}{\sqrt {2 \pi}} \int_ {- \infty} ^ {+ \infty} e ^ {- \xi^ {2} / 2} \mathrm{H} _ {n} (x + \xi) d \xi = x ^ {n}
\]

而 VI, p.~278 的公式 (27) 在 \(h = 0\) 时成为

\[
\begin{array}{l} \sqrt {2 \pi} f (x) = \sum_ {m = 0} ^ {n} \left(\int_ {- \infty} ^ {+ \infty} e ^ {- \xi^ {2} / 2} f ^ {(m)} (\xi) d \xi\right) \frac {\mathrm{H} _ {m} (x)}{m !} \\ - \int_ {- \infty} ^ {+ \infty} d \xi \int_ {0} ^ {\xi} \frac {\mathrm{H} _ {n} (x + \eta)}{n !} e ^ {- (\xi + x) ^ {2} / 2} f ^ {(n + 1)} (x + \xi - \eta) d \eta . * \end{array}
\]

